\subsection{The actual unit-arm curvature for the short-two connector}\label{proof:Q5-C163}
\paragraph{Q163.0 Original goal, exact new scope, and fixed directions}

The surrounding original Q5 problem is the following. For every nonempty finite
original tree U having at most five vertices of original degree at least three,
choose all original edge directions solely from U, before the host, such that
for every original finite probability measure mu and every real skew kernel K
with pointwise \(|K|\le 1,\allowbreak \)
\fitdisplay{
Z_D=\mathbb E_{\mu^{V(U)}}\prod_{i\to j}(1+K(x_i,x_j))\le1.
}{}
Original arm lengths, equal multiplicities, all subdivision vertices, all four
independent positive core lengths, the original edge count, arbitrary original
joint-label laws in entropy applications, and noninjective maps are retained.
We prove the half-tree interface below.

% Replace a pure model-definition block only; keep local root/terminal data.
Use the original finite-space model and arm messages of
Section~\ref{sec:odd-common-calculus}.
When the original \(K\)-notation is used, \(K=-B\) and
\(g=K\one=-u\); in either notation
\(\delta_j=\|B^j\one\|_2^2=\|K^j\one\|_2^2\).
All means, norms, adjoints and spectral integrals use the same
original probability measure.
All expectations, adjoints, norms, and finite spectral integrals use the same
original mu. Until the separate $d=0$ case in \mbox{\hyperref[part:Q163.7]{part Q163.7}}, the proof assumes \(d>0.\)

The original half-tree has root residual arms (1, 2, v), \(v\ge 4,\allowbreak \) and terminal
residual arms (2, b), \(b\ge 5.\) Its inner connector has exactly two original edges.
Use the actual, structure-fixed messages
\fitdisplay{
F=\bar f_2W_b,\qquad P=h_2f_v,\qquad H=pP\,T^2F.
}{Q163.0.1}
The root unit arm is outward, giving \(p=f_1.\) The root length-two arm has its
first edge toward the root and its other edge away, giving \(h_2\); its length-v
arm is all outward, giving \(f_v.\) The terminal length-two arm is all inward,
giving \(\bar f_2\); its length-b arm has its first edge outward and remaining edges
toward the terminal, giving \(W_b.\) The inner connector is directed from the
terminal toward the root, giving \(T^2.\) The independent outer connector is
directed from this root toward its receiving center, giving \(T^M\) for every
original integer \(M\ge 1.\)

Author theorem T171.1. For every original \(b\ge 5,\allowbreak  v\ge 4,\allowbreak \) and \(M\ge 1,\allowbreak \) the single
direction (Q163.0.1), including every genuine equal-pair decoration at the two
original vertices, satisfies
\fitdisplay{
\boxed{EH^4\le e^{-25d/4},\qquad
 I_q(T^M H)^3\le e^{-3d/2},\qquad
 I_q(Y)=E[q^{-1/3}|Y|^{8/3}].}
}{Q163.0.2}
The proof actually gives \(EH^4\le \exp(-6929d/1080).\) In particular it pays the
previously requested ordinary-fourth loss 45d/8 and outer loss d/4.

\paragraph{Q163.1 Precisely used inputs and finite notation}

Use the finite-space bounds \eqref{eq:odd-common:inputs}, the spectral
identities \eqref{eq:odd-common:coordinates}, and the finite convolution
\eqref{eq:Q188:conv}. The fourth signed-mean bound is
\eqref{eq:oc:signed-moments}, proved in Section~\ref{part:Q81.2};
it applies to every original real vector, with its signed mean.
Section~\ref{proof:Q5-C084} gives the ordinary, unanchored outlet
\fitdisplay{
 I_q(T^M H)^3\le e^{11d}(\E|H|^4)^2\qquad(M\ge1).
}{Q163.1.2}

For \(G_n=E_n+zD_n\), the finite coordinate ranges and
\(|G_n|\le1\) at \(n\ge3\) are
Lemma~\ref{lem:odd-common:spectral}. Here the root index is at least
two and the terminal index at least three, so \(G_1\) does not occur.
The genuine \(G_2\) square remains \(1+x\).
The following proof uses its own Taylor coefficient and its own
\(\epsilon,\delta\) and joint-norm estimates.

\paragraph{Q163.2 The same actual terminal and a sharper whole-domain Taylor coefficient}\label{part:Q163.2}

Let \(k=b-2\ge 3\) and write
\fitdisplay{
G_k=A+zC,\quad \gamma=\beta_{b-1}=\int C,\quad
 \Gamma=\beta_b+\gamma=\int(A+C).
}{}
Thus A, C are in [1-x, 1], gamma lies in [kappa d, d], and
\(\Gamma\ge2\kappa d\).

To verify the actual envelope, apply A to the \(\bar f_{b-1}\) convolution while
retaining its original \(\beta_{b-1} \bar f_0\) term. Since \(A1=p\) and
\(E(1+G_{b-1}(B)u)=1-\beta_b\),
\fitdisplay{
W_b\le1-\beta_b-u-BG_{b-1}(B)u-\gamma p.
}{}
As \(\bar f_2=1+u+Bu\) and \(G_{b-1}=1+zG_k\), their sum satisfies
\fitdisplay{
\bar f_2+W_b\le2-\Gamma+R,\qquad
 R=\gamma u-B^2G_k(B)u=(\gamma+xA+zxC)u.
}{}
The two messages are genuine and nonnegative. AM-GM and \(1+t\le e^t\) therefore
give
\fitdisplay{
\boxed{F=\bar f_2W_b\le e^{R-\Gamma}.}
}{Q163.2.1}

The improvement specific to \(b\ge 5\) is the actual finite contraction
\(|G_k|\le1\). The original row bound yields
\fitdisplay{
|R|\le d+\|BG_k(B)u\|\le d+\sqrt{\delta_2}
 \le1/3+\sqrt{13/96}<1/3+3/8=17/24<3/4.
}{Q163.2.2}
The strict square comparison is \(13/96<9/64.\)

For \(0\le t\le 3/4,\allowbreak \) the positive Taylor series gives
\fitdisplay{
{e^t-1-t\over t^2}
 \le {1\over2}+{t\over6}+{t^2\over24(1-t/5)}
 \le {1\over2}+{1\over8}+{15\over544}
 ={355\over544}<{2\over3}.
}{Q163.2.3}
Here the \(n\ge 4\) remainder is bounded using its successive ratio at most t/5.
The last rational comparison is \(1065<1088.\) At \(t=0\) use continuity.
For \(t\le 0,\allowbreak \) the integral Taylor remainder is at most \(t^2/2.\) Hence the same
actual R, everywhere on every original host, obeys
\fitdisplay{
e^R\le1+R+(2/3)R^2.
}{Q163.2.4}
This is \(\eta=4/3,\allowbreak \) not C160's \(\eta=3/2.\) Only an analysis bound changes;
the original terminal F and all its edge directions remain the same.

\paragraph{Q163.3 Actual moments, epsilon loss, and the real constant}

For this same R define
\fitdisplay{
\sigma=ER^2,\quad D_j=EB^jR,\quad
 \delta=D_0+(4/3)\sigma,\quad\epsilon=\delta+D_1.
}{}
The exact finite identities are
\fitdisplay{
\begin{aligned}
 \sigma&=\gamma^2d+2\gamma\int xA
                  +\int x^2(A^2+xC^2),\\
 D_0&=-\int xC,\qquad
 D_1=-\gamma d-\int xA,\qquad
 D_2=\int x^2C.
 \end{aligned}
}{Q163.3.1}
Since \(|G_k|^2=A^2+xC^2\le1\),
\fitdisplay{
\sigma\le\gamma^2d+2\gamma\int xA+\delta_3.
}{Q163.3.2}
No gamma, delta, or sigma is being substituted from a different vector.

\paragraph{Q163.3.1 A uniform negative epsilon}

Let \(X_A= \int xA\), \(X_C= \int xC\).
Using (Q163.3.1)--(Q163.3.2),
\fitdisplay{
\epsilon\le d[-\gamma+(4/3)\gamma^2]
       +[(8/3)\gamma-1]X_A-X_C+(4/3)\delta_3.
}{}
The first term is nonpositive. Since \(\gamma\le 1/3,\allowbreak \) the coefficient of \(X_A\)
is at most -\(1/9.\) Also \(X_A,\allowbreak  X_C\ge \kappa \delta_2\) and \(\delta_3\le \theta \delta_2.\) Therefore
\fitdisplay{
\boxed{\epsilon\le
  -\left({10\kappa\over9}-{4\theta\over3}\right)\delta_2
  =-{17\over144}\delta_2\le0.}
}{Q163.3.3}
This is valid on the entire host domain, without a split at \(d=1/8.\)

\paragraph{Q163.3.2 A sufficient delta bound with a short continuous proof}

The lower bound is immediate: \(\delta\ge -X_C\ge -\delta_2.\) For the upper bound,
\fitdisplay{
\delta\le {4\over3}d^3+
       \left({8\over3}d-\kappa+{4\theta\over3}\right)\delta_2.
}{}
The parenthesized coefficient is first bounded from above by the positive
constant \(241/288\); only then is \(\delta_2\le \theta\) d used.  Since \(d^2\le 1/9,\allowbreak \)
\fitdisplay{
{\delta\over d}\le {4\over27}+{241\theta\over288}
 ={4\over27}+{3133\over9216}<{3\over20}+{7\over20}={1\over2}.
}{Q163.3.4}
The two last comparisons are \(80<81\) and \(62660<64512.\)
Consequently
\fitdisplay{
\boxed{-\delta_2\le\delta<d/2.}
}{Q163.3.5}

\paragraph{Q163.4 Exact two-edge propagation and the entire original root pair}

The positive T kernel is at most two, so \(T(R^2)\le2\sigma\).
Applying T to (Q163.2.1) and then using (Q163.2.4),
\fitdisplay{
TF\le e^{-\Gamma}(q+\delta+BR).
}{}
A second actual application of T gives
\fitdisplay{
T^2F\le e^{-\Gamma}Y,\qquad
 Y=1+u+\epsilon+B V_2,\qquad V_2=\delta\one+u+BR.
}{Q163.4.1}
Indeed \(T\bar f_{1}=\bar f_2=1+u+Bu\) and \(TBR=D_1+B^2R\).
The right side is strictly positive on every positive-mass atom because
it dominates the actual positive message \(T^2F.\) In particular log Y is legal.

Let \(m=v-2\ge 2,\allowbreak  G_m=E+zD\), and
\fitdisplay{
\beta=\beta_{v-1}=\int D,\qquad \beta_v=\int E.
}{}
For the root, retain the actual \(\beta_v f_0\) and \(\beta_{v-1} f_1\) terms in
the positive convolution for \(f_v.\) Since \(h_2=1+u-Bu\) and
\(G_v(-B)=1-B+B^2G_m(-B)\), one obtains
\fitdisplay{
h_2+f_v\le2-\beta_v-\beta
             -B^2G_m(-B)u+\beta u.
}{}
Thus the entire actual root pair obeys
\fitdisplay{
\boxed{P=h_2f_v\le
  \exp\{-\beta_v-\beta+B V_{\rm root}\},\quad
  V_{\rm root}=-BG_m(-B)u+\beta\one.}
}{Q163.4.2}
The actual root unit p remains separate.

Its true scalar-plus-mean loss is
\fitdisplay{
\ell_{\rm root}=\beta_v+\beta-\int xD.
}{}
If m is even then \(E=D,\allowbreak \) so
\fitdisplay{
\ell_{\rm root}=\int(2-x)D
   \ge(2-\theta)(1-\theta)d={969\over1024}d.
}{}
If m is odd then \(E=1-xD,\allowbreak \) so
\fitdisplay{
\ell_{\rm root}=d+\int(1-2x)D
   \ge[1+(1-2\theta)(1-\theta)]d
   ={569\over512}d>{969\over1024}d.
}{}
Therefore, for every actual finite \(v\ge 4,\allowbreak \)
\fitdisplay{
\boxed{\ell_{\rm root}\ge(969/1024)d.}
}{Q163.4.3}
The even minimum \(m=2,\allowbreak \) namely \(v=4,\allowbreak \) is included literally.

Combining the vectors in (Q163.4.1)--(Q163.4.2), define the genuine joint polynomial
vector Q and real scalar c by
\fitdisplay{
Q=[1-xD-x^2C+z(\gamma+xA-E)]u,\qquad
 c=\delta+\beta,\qquad V_2+V_{\rm root}=Q+c1.
}{Q163.4.4}
From (Q163.3.5), \(\beta\ge \kappa\) d, \(\beta\le d,\allowbreak \) and \(\delta_2\le \theta\) d,
\fitdisplay{
\boxed{(3/16)d\le c<(3/2)d\quad(d>0).}
}{Q163.4.5}
In particular c is not zero, and no orthogonality between the constant
vector and Q may be assumed.

The two means are
\fitdisplay{
EQ=\int E-\gamma d-\int xA,\qquad
 E(BQ)=-d+\int xD+\int x^2C.
}{Q163.4.6}
Since \(Eu=0,\allowbreak  E[B(Q+c1)]=E(BQ)\), even though \(EQ\) need not vanish.
The full squared norm remains exactly
\fitdisplay{
\|Q+c1\|^2=\|Q\|^2+c^2+2c\,EQ.
}{Q163.4.7}
The next section pays the pointwise contribution c B1 in full before
applying the exponential lemma.

\paragraph{Q163.5 The original unit arm pays the constant contribution}

Apply \eqref{eq:oc:unit-log} to the actual \(p\).
Using \(\log Y\le Y-1\) in (Q163.4.1)--(Q163.4.2) and the exact
vector (Q163.4.4), set \(a=-\Gamma+\epsilon-\beta_v-\beta\). Then
\fitdisplay{
 \log H\le a+BQ+cu-(3/10)u^2.
}{Q163.5.3}
Lemma~\ref{lem:odd-common:fourth}, with its real coefficient \(c\), gives
\fitdisplay{
 \boxed{\log\E H^4\le4S+(17/4)\|Q\|^2+(10/3)c^2,}
 \qquad S=a+\E(BQ).
}{Q163.5.5}
The pointwise square \eqref{eq:oc:unit-square} pays the entire
\(cu\) term before the exponential estimate is applied to \(Q\).
This keeps the signed mean from (Q163.4.6),
\[
 S=-\Gamma+\epsilon-d+D_2-\ell_{\rm root}.
\]
Thus (Q163.3.3), \(D_2\le\delta_3\le\theta^2d\),
(Q163.4.3), and \(\Gamma\ge2\kappa d\) give
\fitdisplay{
 \boxed{S\le-{95\over32}d-{17\over144}\delta_2.}
}{Q163.5.6}
Here \(2\kappa+1-\theta^2+969/1024=95/32\).

\paragraph{Q163.6 A strengthened joint norm, on the full continuous frequency interval}

Write the two coordinates of Q as
\fitdisplay{
R_0=1-xD-x^2C,\qquad I_0=\gamma+xA-E.
}{}
The actual finite-coordinate ranges imply
\fitdisplay{
0<1-x-x^2\le R_0\le1-x+x^3.
}{}
Also \(I_0\ge x(1-x)-1=-(1-x+x^2)\). On the other side,
\(I_0\le1/3+2x-1\le7/48\), whereas \(1-x+x^2\ge3/4\).
Consequently
\fitdisplay{
|I_0|\le1-x+x^2.
}{}
These inequalities use \(\gamma\le 1/3\) and are valid on the entire original
domain, not just on \(d\le 1/8.\) Thus
\fitdisplay{
\|Q\|^2\le\int
   [1-x-x^2+5x^3-4x^4+x^5+x^6].
}{Q163.6.1}

The following single continuous polynomial estimate pays a useful loss:
\fitdisplay{
1-x-x^2+5x^3-4x^4+x^5+x^6\le1-(4/5)x
 \quad(0\le x\le13/32).
}{Q163.6.2}
After subtracting the right side and dividing by \(x>0,\allowbreak \) it is enough to
check
\fitdisplay{
h(x)=-1/5-x+5x^2-4x^3+x^4+x^5\le0.
}{}
Now
\fitdisplay{
h''(x)=10-24x+12x^2+20x^3\ge10-24\theta=1/4>0.
}{}
A convex function is bounded above by its endpoint chord. The two actual
endpoint values are
\fitdisplay{
h(0)=-1/5,\qquad h(13/32)
       =-{1835407\over167772160}<0.
}{}
This proves (Q163.6.2) throughout the interval; \(x=0\) follows directly.
Therefore
\fitdisplay{
\boxed{\|Q\|^2\le d-(4/5)\delta_2.}
}{Q163.6.3}
The vector Q contains the root and terminal polynomials together.

\paragraph{Q163.7 Complete scalar payment and all outer lengths}\label{part:Q163.7}

Insert (Q163.4.5), (Q163.5.6), and (Q163.6.3) into (Q163.5.5):
\fitdisplay{
\begin{aligned}
 \log EH^4
 &\le-{95\over8}d-{17\over36}\delta_2
      +{17\over4}d-{17\over5}\delta_2+{15\over2}d^2\\
 &=-{61\over8}d-{697\over180}\delta_2+{15\over2}d^2\\
 &\le-{61\over8}d+{653\over180}d^2\\
 &\le-{6929\over1080}d
 <-{25\over4}d\quad(d>0).
 \end{aligned}
}{Q163.7.1}
The last two steps use only the original moment inequality \(\delta_2\ge d^2\)
and the original \(d\le 1/3.\) The exact final excess over \(25/4\) is
\(179/1080>0.\) Hence the ordinary-fourth claim in (Q163.0.2) is proved.

By the exact C084 interface (Q163.1.2), for every independent original
outer length \(M\ge 1,\allowbreak \)
\fitdisplay{
I_q(T^M H)^3
 \le e^{11d}(EH^4)^2
 \le e^{(11-25/2)d}=e^{-3d/2}.
}{Q163.7.2}
This is the ordinary-fourth-to-Iq outlet.

At \(d=0,\allowbreak  u=B\one=0\) on every positive-mass atom. Original leaf integration
makes all actual tree messages equal to one, even when B is not zero.
Thus \(H=1,\allowbreak  q=1\) and both (Q163.0.2) inequalities hold with equality.
No part of the direction changes at that boundary.
